\documentclass[10pt, aspectratio=169]{beamer}

\usepackage[spanish]{babel}
\usepackage[utf8]{inputenc}
\usepackage[pscoord]{eso-pic}
\usepackage{hyperref}
\usepackage{eqnarray,amsmath}
\usepackage[absolute,overlay]{textpos}
\hypersetup{colorlinks=true,urlcolor=blue}

\usepackage{tikz}
\usetikzlibrary{automata, positioning, arrows.meta, petri}

\usepackage{listings}
\usepackage{xcolor}

\usepackage[style=apa]{biblatex}
\bibliography{biblio.bib}

\title{}
\author{}
\date{}

\newtheorem{teorema}{Teorema.}
\setbeamersize{text margin left=10pt,text margin right=10pt}
\beamertemplatenavigationsymbolsempty

\hypersetup{
	colorlinks=true,
	pdfborderstyle={/S/U/W 1}
}

\setbeamertemplate{background}{\includegraphics[width=\paperwidth,height=\paperheight,keepaspectratio]{portada.pdf}}

\newcommand{\azulULA}{\color[rgb]{0.30078125,0.30078125,0.999999}}
\definecolor{azulULA}{RGB}{77, 77, 255}

\setbeamercolor{title}{fg=azulULA}
\setbeamercolor{frametitle}{fg=azulULA}
\setbeamercolor{block title}{bg=azulULA,fg=white}
\setbeamercolor{block body}{bg=azulULA!10!bg}

\renewcommand{\qedsymbol}{\textcolor{black}{\openbox}}

\newcommand{\placetextbox}[3]{
	\setbox0=\hbox{#3}
	\AddToShipoutPictureFG*{
		\put(\LenToUnit{#1\paperwidth},\LenToUnit{#2\paperheight}){\vtop{{\null}\makebox[0pt][c]{#3}}}
	}
}
\newcommand{\placetextboxleft}[3]{
	\setbox0=\hbox{#3}
	\AddToShipoutPictureFG*{
		\put(\LenToUnit{#1\paperwidth},\LenToUnit{#2\paperheight}){\vtop{{\null}\makebox[0pt][l]{#3}}}
	}
}

\begin{document}
	% Diapositiva de la portada.
	\begin{frame}
		\placetextbox{0.47}{0.40}{\large\azulULA\textbf{\begin{tabular}{c} Teoría de la Computación y Lenguajes Formales: \\
					Sesión 11-B. Subclases, Fuentes, Temporización y Workflows \end{tabular}}}
	\end{frame}
	
	% Cambio de formato para las diapositivas centrales (cuerpo/contenido).
	\setbeamertemplate{background}{
		\includegraphics[width=\paperwidth,height=\paperheight,keepaspectratio]{fondo1280x720.pdf}}
	
	\begin{frame}{Alcance de la sesión 11-B}
		\tableofcontents
	\end{frame}

% =======================================================
\section{Subclases Estructurales de Redes de Petri}
% =======================================================

\begin{frame}{}
	\centering
	\Huge{Subclases Estructurales: \\ Grafos de Estados, Grafos de Eventos y Libre Escogencia}
\end{frame}

\begin{frame}{Grafo de Estados (Máquina de Estados)}
	\begin{block}{Definición Formal}
		Una Red de Petri es una \textbf{Máquina de Estados (State Machine)} si cada transición tiene exactamente un lugar de entrada y un lugar de salida:
		\[ \forall t \in T: \quad |{}^\bullet t| = 1 \quad \text{y} \quad |t^\bullet| = 1 \]
	\end{block}

	\begin{columns}
		\column{.55\textwidth}
		\begin{itemize}
			\item \textbf{Capacidad de Modelado:} Modela bifurcaciones y decisiones (conflicto / OR-split), pero \alert{no puede modelar concurrencia (AND-split) ni sincronización (AND-join)}.
			\item \textbf{Conservación de Fichas:} Si $M_0$ tiene una sola ficha, la red conservará exactamente esa ficha en todo momento.
			\item \textbf{Conexión Formal:} Es \textbf{isomorfo a un Autómata Finito (DFA/NFA)}. Cada lugar es un estado del autómata.
		\end{itemize}
		\column{.40\textwidth}
		\begin{center}
			\begin{tikzpicture}[node distance=1.4cm, >=Stealth, thick]
				\node[place, tokens=1] (p1) [label=above:\scriptsize $p_1$] {};
				\node[transition] (t1) [below left of=p1, label=left:\scriptsize $t_1$] {} edge [pre] (p1);
				\node[transition] (t2) [below right of=p1, label=right:\scriptsize $t_2$] {} edge [pre] (p1);
				\node[place] (p2) [below of=t1, label=below:\scriptsize $p_2$] {} edge [pre] (t1);
				\node[place] (p3) [below of=t2, label=below:\scriptsize $p_3$] {} edge [pre] (t2);
			\end{tikzpicture}
		\end{center}
	\end{columns}
\end{frame}

\begin{frame}{Grafo de Eventos (Grafo Marcado / Marked Graph)}
	\begin{block}{Definición Formal}
		Una Red de Petri es un \textbf{Grafo Marcado (Marked Graph o Event Graph)} si cada lugar tiene exactamente una transición de entrada y una transición de salida:
		\[ \forall p \in P: \quad |{}^\bullet p| = 1 \quad \text{y} \quad |p^\bullet| = 1 \]
	\end{block}

	\begin{columns}
		\column{.55\textwidth}
		\begin{itemize}
			\item \textbf{Capacidad de Modelado:} Modela \textbf{concurrencia pura y sincronización} (paralelismo AND-split y AND-join), pero \alert{no puede modelar conflicto, decisiones ni recursos compartidos}.
			\item \textbf{Propiedad Determinista:} Libre de conflictos. No existen decisiones no deterministas.
			\item \textbf{Aplicación:} Modelado de líneas de ensamblaje estrictas y procesamiento en pipeline.
		\end{itemize}
		\column{.40\textwidth}
		\begin{center}
			\begin{tikzpicture}[node distance=1.4cm, >=Stealth, thick]
				\node[transition] (t1) [label=above:\scriptsize $t_1$] {};
				\node[place, tokens=1] (p1) [below left of=t1, label=left:\scriptsize $p_1$] {} edge [pre] (t1);
				\node[place] (p2) [below right of=t1, label=right:\scriptsize $p_2$] {} edge [pre] (t1);
				\node[transition] (t2) [below of=t1, yshift=-1.1cm, label=below:\scriptsize $t_2$] {}
					edge [pre] (p1) edge [pre] (p2);
			\end{tikzpicture}
		\end{center}
	\end{columns}
\end{frame}

\begin{frame}{Redes de Libre Escogencia (Free-Choice Nets)}
	\begin{block}{Definición Formal}
		Una Red de Petri es de \textbf{Libre Escogencia (Free-Choice)} si cuando dos transiciones comparten un lugar de entrada, ese lugar es el único lugar de entrada de ambas:
		\[ \forall t_1, t_2 \in T: \quad {}^\bullet t_1 \cap {}^\bullet t_2 \neq \emptyset \implies {}^\bullet t_1 = {}^\bullet t_2 = \{p\} \]
	\end{block}

	\begin{columns}
		\column{.58\textwidth}
		\begin{itemize}
			\item \textbf{Significado Operacional:} Si una transición en conflicto está habilitada, la otra \textbf{también lo está}. La elección entre ellas depende únicamente de la ficha en $p$, nunca de condiciones de sincronización externas.
			\item \textbf{Sin Confusión:} Separa limpiamente el \textbf{conflicto} de la \textbf{concurrencia}.
			\item \textbf{Importancia en Workflows:} Gran parte de los teoremas de corrección y verificación eficiente de procesos exigen que el flujo sea de Libre Escogencia.
		\end{itemize}
		\column{.38\textwidth}
		\begin{center}
			\begin{tikzpicture}[node distance=1.3cm, >=Stealth, thick]
				\node[place, tokens=1] (p) [label=above:\scriptsize $p$] {};
				\node[transition] (t1) [below left of=p, label=left:\scriptsize $t_1$] {} edge [pre] (p);
				\node[transition] (t2) [below right of=p, label=right:\scriptsize $t_2$] {} edge [pre] (p);
				\node[place] (p1) [below of=t1, label=below:\scriptsize $p_1$] {} edge [pre] (t1);
				\node[place] (p2) [below of=t2, label=below:\scriptsize $p_2$] {} edge [pre] (t2);
			\end{tikzpicture}
		\end{center}
	\end{columns}
\end{frame}

% =======================================================
\section{Concepto de Fuentes y Sumideros}
% =======================================================

\begin{frame}{}
	\centering
	\Huge{Fuentes y Sumideros: \\ Entradas, Salidas y Ciclos de Vida}
\end{frame}

\begin{frame}{Fuentes y Sumideros en Redes de Petri}
	\begin{block}{Nodos Fuente y Nodos Sumidero}
		\begin{tabular}{|l|c|p{4.5cm}|p{4.8cm}|}
			\hline
			\textbf{Elemento} & \textbf{Condición Formal} & \textbf{Significado en el Sistema} & \textbf{Ejemplo} \\ \hline\hline
			\textbf{Lugar Fuente} & ${}^\bullet p = \emptyset$ & Punto de entrada pasivo del sistema. Almacena insumos o solicitudes iniciales. & \texttt{Bandeja\_Entrada\_Pedidos}, \texttt{Buffer\_Materia\_Prima}. \\ \hline
			\textbf{Lugar Sumidero} & $p^\bullet = \emptyset$ & Punto de salida o destino final. Acumula productos terminados o descartes. & \texttt{Almacen\_Despacho}, \texttt{Archivo\_Auditoria}. \\ \hline
			\textbf{Transición Fuente} & ${}^\bullet t = \emptyset$ & Generador autónomo que produce fichas de forma continua o espontánea. & Llegada de clientes según tasa de arribo, generador de eventos externos. \\ \hline
			\textbf{Transición Sumidero} & $t^\bullet = \emptyset$ & Consumidor de fichas que destruye entidades del sistema sin retorno. & Salida del sistema hacia el cliente, desecho de piezas defectuosas. \\ \hline
		\end{tabular}
	\end{block}
\end{frame}

\begin{frame}{Redes de Flujo de Trabajo: Las WF-Nets (van der Aalst)}
	\begin{block}{El Rol Central de Fuente y Sumidero en Workflows}
		Una Red de Petri $PN = (P, T, F)$ es una \textbf{WF-Net} si modela el ciclo de vida completo de un caso:
		\begin{enumerate}
			\item \textbf{Lugar Fuente Único ($i \in P$):} ${}^\bullet i = \emptyset$ (Creación de la instancia del proceso).
			\item \textbf{Lugar Sumidero Único ($o \in P$):} $o^\bullet = \emptyset$ (Culminación de la instancia del proceso).
			\item \textbf{Conectividad Fuerte de Caminos:} Todo nodo $x \in P \cup T$ está en algún camino dirigido desde $i$ hasta $o$ (no hay tareas muertas ni tareas inaccesibles).
		\end{enumerate}
	\end{block}

	\begin{alertblock}{Propiedad de Corrección (\emph{Soundness})}
		Una WF-Net es correcta (\emph{Sound}) si:
		\begin{itemize}
			\item Siempre es posible llegar al estado final $M_f = [o]$ desde cualquier estado alcanzable.
			\item Al alcanzarse el sumidero $o$, \textbf{ningún otro lugar de la red contiene fichas} (limpieza de contexto).
		\end{itemize}
	\end{alertblock}
\end{frame}

% =======================================================
\section{Temporización en Redes de Petri}
% =======================================================

\begin{frame}{}
	\centering
	\Huge{Temporización en Redes de Petri: \\ Del Tiempo Lógico al Tiempo Métrico}
\end{frame}

\begin{frame}{¿Por qué Temporizar una Red de Petri?}
	\begin{block}{La Necesidad en la Ingeniería de Procesos}
		\begin{itemize}
			\item Las Redes de Petri clásicas son **asíncronas**: modelan el orden causal de eventos, no su duración.
			\item En la realidad de la planta y del negocio, las actividades tienen una **duración física** (segundos, horas) y los buffers tienen **tiempos mínimos o máximos de permanencia**.
			\item La temporización permite calcular: \textbf{tiempos de ciclo (lead time), cuellos de botella, productividad y utilización de recursos}.
		\end{itemize}
	\end{block}

	\begin{alertblock}{Las Dos Opciones de Temporización}
		¿Dónde asignamos el retardo temporal $\tau$?
		\begin{enumerate}
			\item \textbf{Asociado a la Transición (T-timed):} La acción o tarea consume tiempo.
			\item \textbf{Asociado al Lugar (P-timed):} La entidad debe reposar o madurar en un estado durante un tiempo.
		\end{enumerate}
	\end{alertblock}
\end{frame}

\begin{frame}{Temporización en la Transición vs. en el Lugar}
	\begin{columns}
		\column{.48\textwidth}
		\begin{block}{1. Temporización en la Transición (T-timed)}
			\begin{itemize}
				\item A cada transición $t$ se le asocia un retardo $\tau(t)$.
				\item \textbf{Semántica:} Al dispararse $t$, consume las fichas de entrada y durante el intervalo $\tau(t)$ la tarea se ejecuta. Al expirar $\tau(t)$, se depositan las fichas en la salida.
				\item \textbf{En Workflows:} Modela la \textbf{duración de una actividad activa} (ej: tiempo de mecanizado CNC, tiempo de redacción de informe por un analista).
			\end{itemize}
		\end{block}

		\column{.48\textwidth}
		\begin{block}{2. Temporización en el Lugar (P-timed)}
			\begin{itemize}
				\item A cada lugar $p$ se le asocia un tiempo de permanencia $\tau(p)$.
				\item \textbf{Semántica:} Cuando una ficha entra en $p$, queda en estado ``latente'' y no puede habilitar transiciones de salida hasta que transcurran $\tau(p)$ unidades de tiempo.
				\item \textbf{En Workflows:} Modela \textbf{tiempos de espera pasivos o maduración} (ej: secado de pintura, enfriamiento de piezas, plazo legal de apelación).
			\end{itemize}
		\end{block}
	\end{columns}
\end{frame}

% =======================================================
\section{Tríptica Holónica, Redes Etiquetadas y Fénix}
% =======================================================

\begin{frame}{}
	\centering
	\Huge{La Tríptica Holónica, \\ Redes Etiquetadas y el Motor Fénix}
\end{frame}

\begin{frame}{Mapeo Holónico: Lugares y Transiciones}
	\begin{block}{Mapeo Operacional de los Lugares ($P$)}
		\begin{itemize}
			\item \textbf{Disponibilidad de Recursos:} Máquina o persona libre para atender trabajo (\texttt{Robot\_Libre}).
			\item \textbf{Etapa de Proceso:} Fase activa de transformación (\texttt{Pieza\_En\_Mecanizado}).
			\item \textbf{Almacén / Buffer:} Capacidad de acumulación de material (\texttt{Buffer\_Entrada\_Horno}).
			\item \textbf{Semáforo Lógico:} Condición que autoriza la ejecución concurrente (\texttt{Autorizacion\_OK}).
		\end{itemize}
	\end{block}

	\begin{block}{Mapeo Operacional de las Transiciones ($T$)}
		\begin{itemize}
			\item \textbf{Tarea Operacional:} Ejecución de una actividad de valor añadido (inicia/termina).
			\item \textbf{Bifurcación (OR-Split):} Decisión excluyente basada en condiciones (aprobar / rechazar).
			\item \textbf{Concurrencia (AND-Split) y Sincronización (AND-Join):} Lanzamiento y consolidación de tareas paralelas.
		\end{itemize}
	\end{block}
\end{frame}

\begin{frame}{Redes Etiquetadas: Triggers y Transiciones Silenciosas ($\epsilon$)}
	\begin{block}{Definición Formal: Función de Etiquetado}
		Dada una Red de Petri y un alfabeto de eventos o \emph{triggers} $\Sigma$, se define:
		\[ \ell: T \longrightarrow \Sigma \cup \{\epsilon\} \]
		\vspace{-0.2cm}
		\begin{itemize}
			\item \textbf{El Trigger ($e \in \Sigma$):} Señal observable por sensor o usuario (ej: \texttt{FIN\_CORTE}, \texttt{PAGO\_OK}).
			\item \textbf{Transición Silenciosa ($\epsilon$):} Acción interna sin sensor; la red evoluciona de $M$ a $M'$ sin evento visible.
		\end{itemize}
	\end{block}

	\begin{columns}
		\column{.48\textwidth}
		\begin{exampleblock}{Perspectiva del Operador}
			\small
			\begin{itemize}
				\item Trabaja en el lenguaje del negocio o planta.
				\item \textbf{No ve $t_{14}$ o $t_{27}$}; refiere el avance de estado exclusivamente por el \textbf{trigger}.
			\end{itemize}
		\end{exampleblock}

		\column{.48\textwidth}
		\begin{alertblock}{Perspectiva del Motor Fénix}
			\small
			\begin{itemize}
				\item Recibe el trigger del sensor o interfaz UI.
				\item Usa el estado actual ($M_k$) para hallar la transición unívoca y avanzar el gemelo digital.
			\end{itemize}
		\end{alertblock}
	\end{columns}
\end{frame}

\begin{frame}{Desambiguación en Fénix: Determinismo por Marcado}
	\begin{block}{El Problema: Sensores de Fin Compartidos en Múltiples Caminos}
		¿Qué ocurre si un sensor señala el mismo fin (\emph{trigger} $e$), pero la pieza o trámite puede provenir de dos caminos distintos (ej: ruta estándar vs. retrabajo)?
	\end{block}

	\begin{columns}
		\column{.46\textwidth}
		\begin{center}
			\begin{tikzpicture}[node distance=1.3cm, >=Stealth, thick]
				\node[place, tokens=1] (p_std) [label={left:\scriptsize $p_{\text{std}}$}] {};
				\node[transition] (t_std) [right of=p_std, label={above:\scriptsize $t_1: \text{\textbf{FIN}}$}] {} edge [pre] (p_std);
				
				\node[place] (p_ret) [below of=p_std, yshift=-0.4cm, label={left:\scriptsize $p_{\text{ret}}$}] {};
				\node[transition] (t_ret) [right of=p_ret, label={below:\scriptsize $t_2: \text{\textbf{FIN}}$}] {} edge [pre] (p_ret);
				
				\node[place] (p_ok) [below right of=t_std, yshift=0.2cm, label={right:\scriptsize $p_{\text{listo}}$}] {};
				\path (t_std) edge [post] (p_ok);
				\path (t_ret) edge [post] (p_ok);
			\end{tikzpicture}
		\end{center}
		{\scriptsize \textbf{Mismo trigger \texttt{FIN}}, pero solo una transición está activa según el camino tomado.}

		\column{.52\textwidth}
		\begin{alertblock}{Determinismo Condicionado al Marcado}
			Con un seguimiento riguroso del estado, para un marcado dado $M_k$, a lo sumo una transición con ese trigger está habilitada:
			\[ |\{ t \in T \mid M_k \ge {}^\bullet t \;\land\; \ell(t) = e \}| \le 1 \]
			\textbf{¡Cero indeterminismo en ejecución!} $M_k$ desambigua el evento.
		\end{alertblock}

		\begin{exampleblock}{Autodiagnóstico en Fénix}
			Si llega el trigger $e$ y el conjunto es vacío ($=0$), Fénix detecta al instante una \textbf{anomalía o evento fuera de secuencia}.
		\end{exampleblock}
	\end{columns}
\end{frame}


\begin{frame}{El Token Coloreado: El Producto con Datos (CPN)}
	\begin{block}{De Fichas Anónimas a Fichas con Atributos}
		En las Redes de Petri Coloreadas (\emph{Coloured Petri Nets - CPN}), las fichas son **objetos de datos tipados**:
		\begin{itemize}
			\item \textbf{El Producto Físico o Documento:} El chasis, la pieza o la orden de compra en tránsito.
			\item \textbf{Atributos Dinámicos:} Cada token porta variables:
			\[ \text{token} = \{ \text{id}: \text{'TK-101'}, \, \text{tipo}: \text{'Acero'}, \, \text{estado}: \text{'Desbastado'} \} \]
			\item \textbf{Funciones de Guarda en las Transiciones:} Una transición solo se dispara si los atributos del token cumplen una condición lógica (ej: \texttt{[prioridad == 'Urgente']}).
		\end{itemize}
	\end{block}

	\begin{alertblock}{La Gran Síntesis}
		$$\text{Lugares (Recursos)} \quad + \quad \text{Transiciones (Procesos)} \quad + \quad \text{Tokens (Productos con Datos)}$$
		Toda la ontología de nuestro sistema se describe bajo un mismo modelo homogéneo y computable.
	\end{alertblock}
\end{frame}

% =======================================================
\section{Casos de Estudio Integrados}
% =======================================================

\begin{frame}{}
	\centering
	\Huge{Casos de Estudio Integrados: \\ Manufactura y Negocios}
\end{frame}

\begin{frame}{Caso 1: Ensamble con Recursos y Temporización}
	\begin{block}{Línea de Ensamble con Sincronización y Tiempos de Maquinado}
		Dos líneas paralelas preparan Chasis y Motor. Se ensamblan mediante un Robot de Montaje ($t_{\text{montar}}$ con $\tau = 12\text{s}$):
	\end{block}

	\begin{center}
		\begin{tikzpicture}[node distance=2.2cm, >=Stealth, thick]
			\node[place, tokens=1] (chasis) [label={above:\scriptsize Chasis Listo}] {};
			\node[place, tokens=1] (motor) [below of=chasis, yshift=-0.5cm, label={below:\scriptsize Motor Listo}] {};
			\node[place, tokens=1] (robot) [right of=chasis, yshift=-1.35cm, label={above:\scriptsize Robot Libre}] {};
			
			\node[transition] (montar) [right of=robot, label={above:\scriptsize $t_{\text{montar}} (\tau=12\text{s})$}] {}
				edge [pre] (chasis)
				edge [pre] (motor)
				edge [pre] (robot);
				
			\node[place] (ensamblado) [right of=montar, label={above:\scriptsize Vehículo Ensamblado}] {}
				edge [pre] (montar);
			\path (montar) edge [post, bend left=45] (robot);
		\end{tikzpicture}
	\end{center}
	\centering{\footnotesize Sincronización AND-Join estricta: requiere Chasis + Motor + Robot simultáneamente.}
\end{frame}

\begin{frame}{Caso 2: Proceso de Compras entre 4 Dependencias}
	\begin{block}{Fuente, Sumideros y Bifurcaciones que Cortan el Flujo}
		El pedido inicia en el \textbf{lugar fuente} $p_{\text{sol}}$ y culmina en el \textbf{lugar sumidero} $p_{\text{fin}}$ o aborta en $p_{\text{rech}}$:
	\end{block}

	\begin{center}
		\begin{tikzpicture}[node distance=2.6cm, >=Stealth, thick]
			\node[place, tokens=1] (p_sol) [label={above:\tiny $i$ (Solicitud)}] {};
			\node[transition] (t_sol) [right of=p_sol, label={above:\tiny solicitar}] {} edge [pre] (p_sol);
			\node[place] (p_eval) [right of=t_sol, label={above:\tiny Evaluando}] {} edge [pre] (t_sol);
			
			\node[transition] (t_ord) [right of=p_eval, label={above:\tiny gen\_orden}] {} edge [pre] (p_eval);
			\node[place] (p_rec) [right of=t_ord, label={above:\tiny Recepción}] {} edge [pre] (t_ord);
			
			\node[transition] (t_pag) [right of=p_rec, label={above:\tiny pagar}] {} edge [pre] (p_rec);
			\node[place] (p_fin) [right of=t_pag, label={above:\tiny $o$ (Pagado)}] {} edge [pre] (t_pag);

			% Bifurcaciones de rechazo
			\node[place] (p_rech) [below of=t_ord, yshift=-0.8cm, label={below:\tiny $o_{\text{abort}}$ (Rechazo)}] {};
			\node[transition] (t_r1) [below of=p_eval, label={left:\tiny rech\_pedido}] {} edge [pre] (p_eval) edge [post] (p_rech);
			\node[transition] (t_r2) [below of=p_rec, label={right:\tiny rech\_prod}] {} edge [pre] (p_rec) edge [post] (p_rech);
		\end{tikzpicture}
	\end{center}
\end{frame}

% =======================================================
\section{Conclusión}
% =======================================================

\begin{frame}{Conclusión de la Sesión 11-B}
	\begin{block}{Lo que aprendimos en esta sesión}
		\small
		\begin{itemize}
			\item \textbf{Subclases:} Grafo de Estados (autómata/conflicto), Grafo de Eventos (concurrencia pura) y Libre Escogencia (sin confusión).
			\item \textbf{Fuentes y Sumideros:} Entradas y salidas del ciclo de vida en WF-Nets con \emph{Soundness}.
			\item \textbf{Temporización:} T-timed (duración activa de tareas) vs. P-timed (maduración/espera en lugares).
			\item \textbf{La Tríptica Unificada:} Recursos (Lugares), Procesos (Transiciones) y Productos (Tokens CPN).
			\item \textbf{Redes Etiquetadas y Motor Fénix:} Triggers del operador vs. transiciones internas ($t$). Desambiguación unívoca por marcado $M_k$ (cero indeterminismo en ejecución) y detección de anomalías.
		\end{itemize}
	\end{block}

	\begin{exampleblock}{Próximo Bloque (Semana 12): Análisis Formal y Propiedades}
		\small
		Verificación formal: **Árbol de Alcanzabilidad, Invariantes de Lugar ($P$-invariantes) y Detección de Bloqueos (\emph{Deadlocks})**.
	\end{exampleblock}
\end{frame}

% Cambio de formato para la diapositiva final.
\setbeamertemplate{headline}{}
\setbeamertemplate{background}
{\includegraphics[width=\paperwidth,height=\paperheight,keepaspectratio]{salida.pdf}}

% Diapositiva final.
\begin{frame}
	\placetextbox{0.6}{0.87}{\Huge\textbf{¡Muchas gracias!}}
\end{frame}

\setbeamertemplate{headline}{}
\setbeamertemplate{background}{}

\printbibliography
\end{document}
